\chapter[Log-Likelihood Field]{Functional Calculus on the Log-Likelihood Field}
\section{Introduction}
The previous chapter revealed that both the Kullback--Leibler divergence
and the Universal Optimality Defect (UOD) are naturally expressed in
terms of the same logarithmic quantity. This chapter develops that
observation into a unified functional framework.
Throughout,
\[ a(P,Q)=\left(a_i\right),\qquad
a_i=\log\frac{P_i}{Q_i}, \]
will be called the \textbf{log-likelihood field}.
\section{The Log-Likelihood Field}
The logarithmic embedding satisfies
$L(P)-L(Q)=a.$
Consequently every quantity introduced in Part V is obtained by
applying a functional to the same geometric object.
The field (a) represents the logarithmic displacement between two
probability distributions in quotient coordinates.
\section{First Moment: Kullback--Leibler Divergence}
\begin{proposition}
\label{ch:prop:27-1}
The Kullback--Leibler divergence is the expectation of the
log-likelihood field with respect to the distribution (P):
\[ \boxed{
D_{KL}(P\|Q)=E_P[a].
} \]
\textbf{Proof.}
By definition,
\[ D_{KL}(P\|Q) = \sum_{i}P_i\log\frac{P_i}{Q_i}
= \sum_{i}P_ia_i = E_P[a]. \]
\end{proposition}
\section{Reverse KL Divergence}
\begin{proposition}
\label{ch:prop:27-2}
The reverse Kullback--Leibler divergence satisfies
\[ \boxed{
D_{KL}(Q\|P)
=
-E_Q[a].
} \]
\textbf{Proof.}
Since
$\log\frac{Q_i}{P_i}=-a_i,$
we obtain
\[ D_{KL}(Q\|P) = \sum_{i}Q_i\log\frac{Q_i}{P_i}
= -\sum_{i}Q_ia_i. \]
\end{proposition}
\section{Second Centered Uniform Moment}
Chapter 21 proved that
\[ \boxed{
\mathcal{D}(P,Q)
=
n\,\operatorname{Var}_U(a),
} \]
where
$U=\left(\frac1n,\ldots,\frac1n\right).$
Equivalently,
$\frac{\mathcal{D}}{n} = E_U[a^{2}]-E_U[a]^{2}.$
Thus the UOD is the centered second moment of the same log-likelihood
field, computed with respect to the uniform measure.
\section{A Unified Perspective}
The previous propositions show that several apparently unrelated
quantities are obtained by evaluating different statistics of the same
field.
\[
\begin{array}{lll}
\text{Functional} & \text{Measure} & \text{Observable} \\
\hline
(D_{KL}(P\|Q)) & P & a \\
(D_{KL}(Q\|P)) & Q & -a \\
\text{UOD} & \text{Uniform} & (a-E_U[a])^2
\end{array}
\]
This table separates two independent choices:
\begin{itemize}
  \item the reference measure;
  \item the functional applied to the field.
\end{itemize}
\section{Functional Viewpoint}
The previous examples motivate the definition of a general functional
$F_{\mu,\Phi}(P,Q) = E_\mu[\Phi(a)],$
where
\begin{itemize}
  \item \(\mu\) is a probability measure;
  \item \(\Phi\) is a functional acting on the log-likelihood field.
\end{itemize}
Different choices recover different geometric quantities.
This observation motivates the general theory developed in the next
chapter.
\section{Interpretation}
The log-likelihood field should be regarded as the fundamental geometric
object associated with the logarithmic quotient embedding.
Classical divergences are not primary objects; rather, they arise by
applying different statistical functionals to the same field.
Within this interpretation, the UOD occupies a distinguished position as
the uniform variance of the logarithmic displacement.
\section{Summary}
This chapter established that
\begin{itemize}
  \item the log-likelihood field is the common geometric object underlying
\end{itemize}
    several divergences;
\begin{itemize}
  \item the KL divergence is its expectation under (P);
  \item the reverse KL divergence is its expectation under (Q) with opposite
\end{itemize}
    sign;
\begin{itemize}
  \item the UOD is its centered second moment under the uniform measure.
\end{itemize}
These identities motivate the general functional calculus developed in
the next chapter.
